\chapter{基础知识}

\dirtree{%
.1 A/B Test 四件事.
.2 1. 让两组可比.
.3 随机分流.
.3 对照组 / 实验组.
.3 排除混杂变量.
.3 选择随机化单位.
.3 核心：除了实验改动，其他都尽量一样.
.2 2. 定义要比较什么.
.3 核心指标 / OEC.
.3 业务指标.
.3 诊断指标.
.3 护栏指标.
.3 核心：差异必须落在有业务意义的指标上.
.2 3. 判断差异是否可信.
.3 p-value.
.3 置信区间.
.3 样本量.
.3 power.
.3 方差 / 标准误.
.3 A/A test.
.3 SRM.
.3 核心：这个差异是实验效果，还是随机波动 / 数据问题？.
.2 4. 决定是否上线.
.3 收益是否足够大.
.3 护栏是否受损.
.3 是否需要灰度放量.
.3 是否需要继续实验.
.3 核心：A/B test 的终点是产品决策.
}


\dirtree{%
.1 Trustworthy Online Controlled Experiments：用线上随机实验判断产品/算法/策略改动是否真的有效.
.2 0 基础工具层：不是目的，是后面能力的弹药.
.3 随机化与因果推断.
.4 要点：通过随机分组，让实验组和对照组可比，从而判断因果影响.
.3 实验指标体系.
.4 要点：知道用什么指标判断实验成功，以及用什么指标防止副作用.
.3 统计推断基础.
.4 要点：理解均值、方差、置信区间、p-value、统计功效、样本量.
.3 实验数据基础.
.4 要点：理解日志、埋点、分流、实验单元、指标口径和数据质量.
}
\begin{definition}[综合评估标准]
综合评估标准（Overall Evaluation Criterion, OEC）是实验目标的定量测量。
\end{definition}
\begin{definition}[多臂老虎机]
多臂老虎机（Multi-Armed Bandit, MAB）是一类序贯决策问题：在每一轮中，决策者需要从多个可选方案（arms）中选择一个执行，并观察该方案带来的随机奖励。由于每个方案的真实收益分布未知，决策者需要在\textbf{探索}（尝试不同方案以学习其收益）和\textbf{利用}（选择当前估计收益最高的方案以获得回报）之间进行权衡，目标是在整个实验过程中最大化累计奖励，或等价地最小化累计遗憾（regret）。
\end{definition}
\begin{itemize}
    \item \textbf{A/B 实验}
    
    A/B 实验的核心思想是先公平比较，再做决策。它通常采用固定分流，例如 A 组 50\%、B 组 50\%，在实验期间保持流量比例不变，等样本量足够后再比较不同方案的指标表现。A/B 实验更关注新方案是否显著优于旧方案，因此适合产品改版、推荐算法评估、补贴策略评估等需要严谨因果结论的场景。它的优点是统计推断更干净、可解释性更强；缺点是如果某个版本明显较差，实验期间仍然会分到固定流量，可能造成机会成本。

    \item \textbf{多臂老虎机}

    多臂老虎机的核心思想是一边探索，一边利用当前表现更好的方案。它不会长期保持固定分流，而是根据不同方案的实时表现动态调整流量，把更多用户分配给当前收益更高的版本，同时保留一部分流量继续探索其他版本。多臂老虎机更关注在不确定环境下如何最大化累计收益，因此适合广告素材、推荐位、标题、封面等需要快速优化收益的场景。它的优点是能够减少流量浪费、降低机会成本；缺点是由于分流比例不断变化，统计推断和因果解释会更加复杂，不能简单套用传统 A/B 实验的显著性检验。
\end{itemize}
\begin{dxtips}
    统计里面p值一半都用于否定假设
\end{dxtips}
\begin{dxtips}
    影响选取用户数量多少都考虑，这个核心思想就是控制变量，确保无关变量不要影响实验
\end{dxtips}
\begin{itemize}
    \item \textbf{指标方差：} 指标波动越大，越需要大样本；指标越稳定，所需样本越小。
    \item \textbf{最小可检测效应：} 想检测的变化越小，越需要大样本；只关心较大变化时，可以使用较小样本。
    \item \textbf{显著性水平：} 显著性水平越严格，例如从 $0.05$ 降到 $0.01$，越需要大样本。
    \item \textbf{统计功效：} 功效要求越高，例如从 $80\%$ 提到 $90\%$，越需要大样本。
    \item \textbf{实验风险：} 风险越高的实验，初始放量应越小，先小流量验证再逐步扩大。
    \item \textbf{流量竞争：} 多个实验同时运行会分走流量，使单个实验可用样本变少，通常需要延长实验时间。
    \item \textbf{时间效应：} 若存在周内效应、季节性、节假日或新奇效应，实验需要覆盖足够长的时间，避免短期数据误导结论。
\end{itemize}
\begin{definition} [魏特曼定理]
    数据异常往往是做错了
\end{definition}
\subsection{辛普森悖论}
\begin{example}[辛普森悖论]
比较旧版本 A 和新版本 B 的转化率。用户分为新用户和老用户。

\begin{center}
\begin{tabular}{c|cc}
\hline
用户类型 & A 组 & B 组 \\
\hline
新用户 & $10/100=10\%$ & $110/1000=11\%$ \\
老用户 & $720/900=80\%$ & $81/100=81\%$ \\
\hline
\end{tabular}
\end{center}

分组内看，B 都比 A 好：
\[
11\%>10\%, \qquad 81\%>80\%.
\]

但合并后：
\[
A:\quad \frac{10+720}{100+900}
= \frac{730}{1000}
= 73\%,
\]
\[
B:\quad \frac{110+81}{1000+100}
= \frac{191}{1100}
\approx 17.4\%.
\]

所以总体上反而是
\[
A>B.
\]

原因是 A 组老用户占比高，而 B 组新用户占比高。老用户本来就更容易转化，因此直接合并会让总体平均被用户结构扭曲。

这说明：在 A/B Test 中，如果两组人群结构不平衡，只比较总体指标可能得到错误结论。
\end{example}
\begin{dxtips}
    也就是辛普森悖论要求我们求这个指标的时候用人群加权过后的平均值而不是直接求平均。因为就是老用户转化率都高，但是A组的老客户人数多，B组老客户人数少，所以在新客户转化率相差不大的情况下，A组的转化率总人数就多
\end{dxtips}
\chapter{指标}
\begin{definition}[目标指标]
目标指标也称为成功指标或北极星指标，是组织最终关心的核心结果。
它回答的是：产品为什么存在，组织的成功是什么样的。
目标指标通常与组织使命直接相关，用来刻画长期、全局、最终的业务成功。
例如，用户长期留存、总观看时长、平台交易额、用户满意度等都可以作为目标指标。
\end{definition}

\begin{definition}[驱动指标]
驱动指标是能够影响目标指标的中间指标，也可以理解为目标指标的可操作分解。
它回答的是：哪些用户行为或业务过程会推动最终成功。
驱动指标通常比目标指标更短期、更敏感，因此常用于实验分析和产品迭代。
例如，点击率、完播率、互动率、人均使用时长等都可以作为驱动指标。
\end{definition}

\begin{definition}[护栏指标]
护栏指标是用来防止实验或策略带来负面影响的约束性指标。
它回答的是：在优化目标指标或驱动指标时，哪些底线不能被破坏。
护栏指标不一定要求显著提升，但通常要求不能显著下降。
例如，崩溃率、延迟、投诉率、取消关注率、内容低质率等都可以作为护栏指标。
\end{definition}
\begin{example}[短视频推荐场景]
假设一个短视频平台希望通过推荐系统提升用户体验。

此时，目标指标可以设为用户总观看时长或长期留存率，因为它们反映了平台最终希望提升的核心结果。

驱动指标可以包括点击率、完播率、单次播放时长和互动率，因为这些用户行为会直接影响总观看时长和留存。

护栏指标可以包括投诉率、取消关注率、低质内容曝光率和应用崩溃率，因为即使推荐策略提升了观看时长，也不能以牺牲内容质量、用户体验或系统稳定性为代价。
\end{example}
\begin{definition}[OEC]
OEC（Overall Evaluation Criterion，整体评价准则）是 A/B 实验中用来判断实验方案是否成功的总体评价指标。
它通常由一个或多个核心指标组成，用来把产品目标转化为可量化、可比较的实验判断标准。

在实验决策中，OEC 回答的是：
如果实验组和对照组相比发生了变化，我们应该根据什么标准判断这个变化是否值得上线。
\end{definition}
\begin{dxtips}
    也就是设计很多指标每一个指标都有一定的权重
\end{dxtips}